\chapter{Infraestructura y despliegue}
\label{chap:infraestructura-despliegue}

\section{Contenerización}
\label{sec:contenerizacion}

La API se empaqueta en una imagen Docker construida a partir de su \texttt{Dockerfile}. El contenedor expone el puerto \texttt{8080}, instala las dependencias necesarias para el runtime .NET y se configura mediante variables de entorno. La imagen se utiliza tanto para despliegue cloud como para validar el comportamiento de empaquetado.

El entorno local se reproduce con Docker Compose, levantando la API y PostgreSQL. La configuración sensible se proporciona mediante un fichero \texttt{.env} no versionado. Esto permite ejecutar el sistema local sin instalar PostgreSQL de forma manual y mantiene una ruta de validación cercana al despliegue real.

\section{Infraestructura como código}
\label{sec:infraestructura-codigo}

Terraform se organiza bajo \path{infra/terraform}. La estructura separa un bootstrap de estado remoto, módulos reutilizables y raíces de entorno:

\begin{itemize}
  \item \path{bootstrap/remote_state}: crea el bucket S3 de estado y la tabla DynamoDB de bloqueo.
  \item \path{modules/platform_foundation}: VPC, subredes, rutas, NAT y grupos de seguridad.
  \item \path{modules/container_registry}: repositorio ECR.
  \item \path{modules/database}: RDS PostgreSQL.
  \item \path{modules/runtime_config}: Secrets Manager y SSM.
  \item \path{modules/container_runtime}: ECS, task definition, ALB, listeners, Route 53 y ACM.
  \item \path{modules/observability}: grupo de logs, dashboard CloudWatch, alarmas y notificación SNS opcional.
  \item \path{modules/github_oidc}: rol de despliegue para GitHub Actions.
  \item \path{environments/dev}: composición real del entorno de desarrollo.
\end{itemize}

El estado remoto se almacena en S3 y se bloquea con DynamoDB. Esto evita estados locales divergentes y reduce el riesgo de cambios concurrentes.

\section{Entorno AWS}
\label{sec:entorno-aws}

El entorno desplegado utiliza:

\begin{itemize}
  \item Cuenta AWS: \texttt{661000947340}.
  \item Región: \texttt{eu-west-1}.
  \item Entorno: \texttt{dev}.
  \item Dominio previsto: \path{api.dev.transitops.net}.
  \item Hosted zone: \path{transitops.net} en Route 53.
  \item Identificador de hosted zone: \path{Z0844787W37HXN9FIJR}.
  \item Fallback operativo: DNS público del ALB cuando el entorno se recrea antes de completar DNS/ACM.
\end{itemize}

Los recursos siguen la convención \path{transitops-dev-<componente>} y se etiquetan con proyecto, entorno, propietario, repositorio, servicio, origen Terraform y las etiquetas \texttt{TerraformStack} y \texttt{ResourceGroup}. Esta nomenclatura permite identificar recursos y costes, y comprobar tras un \texttt{terraform destroy} que no quedan recursos del entorno \texttt{dev}.

Como el entorno \texttt{dev} es educativo y desechable, la configuración actual desactiva la retención automática de backups de RDS y omite snapshots finales. Esta decisión reduce costes residuales en ciclos frecuentes de creación y destrucción, sin impedir que un entorno futuro de producción defina una política de backup más conservadora.

\section{Despliegue de la aplicación}
\label{sec:despliegue-aplicacion}

El flujo de despliegue aplicado en \texttt{dev} durante la primera validación cloud fue:

\begin{enumerate}
  \item Crear la infraestructura base con ECS en \texttt{desired\_count} = \texttt{0}, evitando arrancar tareas antes de tener imagen en ECR.
  \item Construir la imagen Docker de la API.
  \item Publicar la imagen en ECR con el tag \path{sprint4-local-d00aa6b}.
  \item Actualizar la task definition para usar la imagen publicada.
  \item Cargar los secretos de runtime en Secrets Manager.
  \item Ejecutar migraciones con una tarea ECS puntual usando \path{--migrate-only}.
  \item Escalar el servicio ECS a \texttt{desired\_count} = \texttt{1}.
  \item Activar Route 53, ACM, HTTPS y redirección HTTP a HTTPS.
  \item Validar el endpoint público de readiness sobre HTTPS.
\end{enumerate}

En Sprint 6 se repitió la recreación del entorno en la cuenta AWS actual \texttt{661000947340}. La imagen validada fue \path{sprint6-local-6bb910c}; se restauraron e importaron secretos que habían quedado programados para eliminación tras un ciclo anterior, se ejecutó una tarea ECS de migración con código de salida \texttt{0}, se escaló el servicio a \texttt{desired = 1} y se validaron readiness, bootstrap de administrador y login. La validación funcional principal de Sprint 6 se hizo mediante el DNS HTTP del ALB, y posteriormente se corrigió la delegación DNS de \path{transitops.net} en la misma cuenta para que el siguiente \texttt{terraform apply} pueda converger directamente con \texttt{enable\_https = true}.

El módulo de runtime incluye además el deployment circuit breaker de ECS con rollback, periodo de gracia para health checks y parámetros explícitos para la comprobación del target group: ruta \path{/api/v1/health/ready}, intervalo, timeout y umbrales de healthy/unhealthy. Esto reduce el riesgo de dejar una versión fallida recibiendo tráfico.

\section{DNS, certificados y HTTPS}
\label{sec:dns-certificados-https}

El endpoint objetivo es \url{https://api.dev.transitops.net}. Para lograrlo, Terraform solicita un certificado público en ACM para \path{api.dev.transitops.net}, crea el CNAME de validación en Route 53, espera a que el certificado pase a estado \texttt{ISSUED}, crea el listener HTTPS \texttt{443} en el ALB y convierte el listener HTTP \texttt{80} en una redirección permanente hacia HTTPS.

Durante la activación en la cuenta \texttt{661000947340} apareció una incidencia real de DNS: el dominio registrado \path{transitops.net} apuntaba a nameservers distintos de los de la hosted zone usada por Terraform. ACM permanecía en \texttt{PENDING\_VALIDATION} aunque el CNAME existía en Route 53, porque el DNS público no consultaba esa zona. La corrección consistió en actualizar los nameservers del dominio registrado para que coincidieran con la hosted zone \path{Z0844787W37HXN9FIJR}. Tras la propagación, el CNAME de validación fue visible desde DNS público y el certificado pasó a \texttt{ISSUED}.

\begin{figure}[htbp]
\centering
\fbox{\begin{minipage}{0.88\textwidth}
\centering
\vspace{1.2cm}
\textbf{Placeholder de captura}\\
Captura de ACM con el certificado de \path{api.dev.transitops.net} en estado \texttt{Issued} y validación DNS correcta.
\vspace{1.2cm}
\end{minipage}}
\caption{Evidencia prevista de certificado ACM emitido.}
\label{fig:placeholder-acm-issued}
\end{figure}

\section{Cierre del entorno y control de costes}
\label{sec:cierre-costes}

Después de capturar la evidencia de Sprint 6, el entorno \texttt{dev} se destruyó con \texttt{terraform destroy} para evitar gasto evitable. La destrucción eliminó ALB, listeners, certificado ACM, registros DNS de \path{api.dev.transitops.net}, ECS, ECR, RDS, NAT Gateway, VPC, subredes, grupos de seguridad, CloudWatch logs, dashboard, alarmas y secretos del entorno. El repositorio ECR necesitó una limpieza previa de imágenes, ya que AWS no permite eliminar un repositorio no vacío.

Tras el destroy se verificó que el state de Terraform quedaba vacío y que no existían recursos costosos asociados al stack \path{transitops-dev}: ALB, RDS, ECS, ECR, certificado ACM, registros DNS del subdominio, log groups, dashboards, alarmas ni secretos. Permanecen fuera del destroy el dominio registrado \path{transitops.net}, la hosted zone pública y los recursos de backend remoto de Terraform, que son recursos base de bajo coste y necesarios para recrear el entorno.

\section{Gestión de configuración y secretos}
\label{sec:configuracion-secretos}

La configuración sensible se externaliza. La cadena de conexión a PostgreSQL, la clave de firma JWT y el token de bootstrap se almacenan en Secrets Manager. Otros parámetros, como emisor, audiencia y expiración del JWT, se almacenan en SSM Parameter Store o se proporcionan como variables no secretas.

En el repositorio solo se incluyen ejemplos y nombres de variables. Los ficheros locales que pueden contener secretos, como \texttt{terraform.tfvars}, \texttt{backend.hcl} o \texttt{.env}, están excluidos del control de versiones.
